\section{Finite identities behind the analytic representations}\label{sec:series}

Throughout, $0^0=1$, $B_0=1$, and a sum over an empty range is zero.  The
Stirling numbers and Bell numbers are the primitive recursive integers defined
in Section~\ref{sec:finite}.  We first prove exact truncation identities; the
infinite notation is then a consequence with an explicit error, not a premise.

\subsection{Dobi\'nski's formula with a uniform arithmetic error}

Write $(x)_j=x(x-1)\cdots(x-j+1)$ and $(x)_0=1$.  The polynomial identity
\[
 x^n=\sum_{j=0}^n S(n,j)(x)_j
\]
follows either from the Stirling recurrence or by classifying functions
$[n]\to[x]$ according to their nonempty fibers.  Substituting the integers
$k=0,\ldots,M$ and interchanging two \emph{finite} sums gives, for $M\ge n$,
\begin{equation}\label{eq:finite-dobinski}
 A_{n,M}:=\sum_{k=0}^M\frac{k^n}{k!}
 =\sum_{j=0}^n S(n,j) E_{M-j}(1).
\end{equation}
The cancellation $(k)_j/k!=1/(k-j)!$ is used only when $k\ge j$; for $k<j$
the falling factorial vanishes.

\begin{theorem}[Dobi\'nski, with a finite certificate]\label{thm:dobinski}
For $M\ge n$,
\begin{equation}\label{eq:dobinski-error}
 0\le B_n-e^{-1}A_{n,M}\le
 \frac{2B_n}{(M-n+1)!}.
\end{equation}
Consequently
\[
 B_n=e^{-1}\sum_{k=0}^{\infty}\frac{k^n}{k!}.
\]
For absolute error at most $2^{-s}$, it is sufficient to take
$M=n+s+(n+1)^2+2$, and to approximate $e^{-1}$ to the additional precision
needed for multiplication by the explicit rational $A_{n,M}$.
\end{theorem}
\begin{proof}
Subtract \eqref{eq:finite-dobinski} from
$eB_n=e\sum_{j=0}^nS(n,j)$.  Each difference $e-E_{M-j}(1)$ is nonnegative
and at most $2/(M-j+1)!\le2/(M-n+1)!$.  Sum the inequalities and divide by
$e\ge1$.  Finally, $B_n\le(n+1)^{n+1}\le2^{(n+1)^2}$ and
$d!\ge2^{d-1}$ for $d\ge1$, both by finite induction.  These inequalities
verify the claimed truncation rule.  If the numerical value of $A_{n,M}$ is
large, choose $t=s+\lceil A_{n,M}\rceil+3$ for the reciprocal's precision;
then $A_{n,M}2^{-t}\le2^{-s-2}$.  Every operation in this rule is primitive
recursive.
\end{proof}

\begin{corollary}[The Poisson interpretation and its polynomial extension]
The weights $w_k=e^{-1}/k!$ have total mass $1$ and their $n$th moment is $B_n$.
For a nonnegative rational parameter $u$, define
$\Phi_n(u)=\sum_{j=0}^n S(n,j)u^j$.  Then
\begin{equation}\label{eq:poisson-touchard}
 e^{-u}\sum_{k=0}^{\infty}\frac{k^n u^k}{k!}=\Phi_n(u),
 \qquad \Phi_{n+1}(u)=u\bigl(\Phi_n(u)+\Phi_n'(u)\bigr).
\end{equation}
\end{corollary}
\begin{proof}
For finite $M\ge n$ the weighted version of
\eqref{eq:finite-dobinski} is
$\sum_{k\le M}k^nu^k/k!=\sum_{j\le n}S(n,j)u^j E_{M-j}(u)$.
Apply \eqref{eq:exp-tail}.  The derivative identity follows by substituting
$S(n+1,j)=jS(n,j)+S(n,j-1)$ and collecting polynomial coefficients.  There is
no construction of a sigma-algebra or arbitrary random variable here:
``moment'' denotes this specified effective weighted sum.
\end{proof}

\subsection{A finite ceiling formula}

\begin{theorem}[Comtet's truncation]\label{thm:comtet}
For every integer $n\ge1$,
\[
 B_n=\left\lceil e^{-1}\sum_{k=1}^{2n}\frac{k^n}{k!}\right\rceil.
\]
\end{theorem}
\begin{proof}
Set $t_k=k^n/k!$.  For $k\ge2n+1$,
\[
 \frac{t_{k+1}}{t_k}
 =\frac{(1+1/k)^n}{k+1}
 \le\frac{e^{n/k}}{k+1}
 <\frac2{k+1}\le\frac12.
\]
Here $e^{1/2}<2$, for example by squaring and using $e<3<4$.
Let $a_n=(2n+1)^n/(2n+1)!$.  Then $a_1=1/2$ and
\[
 \frac{a_{n+1}}{a_n}
 =\frac{(1+2/(2n+1))^n}{2n+2}
 <\frac3{2n+2}\le\frac34.
\]
Thus $a_n\le1/2$, and the omitted positive tail is less than $2a_n\le1$.
By Theorem~\ref{thm:dobinski}, the truncated expression lies strictly between
$B_n-1$ and $B_n$.  This is precisely the ceiling assertion.  The constants
$e<3$ and $(1+x)^m\le e^{mx}$ follow from the factorial series and the binomial
expansion.  This proof does not assume that the ceiling of an arbitrary real
name is computable: it proves a specific interval with an integer endpoint.
\end{proof}

\subsection{Exponential generating functions and complete Bell polynomials}

A formal power series identity below means equality of every finite truncated
coefficient list.  In $\Q[z]/(z^{N+1})$, substitution into an exponential is
finite whenever the substituted polynomial has constant coefficient zero.

\begin{theorem}[The exponential generating function]\label{thm:egf}
As a formal series, and also as an effective complex analytic identity,
\begin{equation}\label{eq:egf}
 \sum_{n=0}^{\infty}\frac{B_n}{n!}z^n=\exp(e^z-1).
\end{equation}
In particular, for every real $r>0$ and every $n\ge0$,
\begin{equation}\label{eq:coefficient-bound}
 B_n\le n!\,r^{-n}\exp(e^r-1).
\end{equation}
\end{theorem}
\begin{proof}
The coefficient of $z^n$ in $(e^z-1)^k/k!$ counts ordered lists of $k$
nonempty labeled blocks, divided by their $k!$ possible orders; it is
$S(n,k)/n!$.  In detail, choose block sizes
$d_1+\cdots+d_k=n$ with $d_i\ge1$.  The multinomial count is
$n!/(d_1!\cdots d_k!)$; dividing by $n!k!$ gives exactly the coefficient
obtained by multiplying the $k$ series.  All $k>n$ contribute zero.
Summing $k=0,\ldots,n$ proves the formal identity.

There is also a direct recurrence check: if $F(z)=\sum B_nz^n/n!$, the
binomial recurrence gives $F'=e^zF$ and $F(0)=1$.  The formal exponential on
the right of \eqref{eq:egf} has the same equations.  Coefficient induction
therefore gives the same result without formal division by an unknown series.

For $r>0$, every finite partial sum of $F(r)$ is bounded above by
$\sum_{k\ge0}(e^r-1)^k/k!=\exp(e^r-1)$: expand the finite coefficients just
computed, and only add nonnegative terms.  The single $n$th term yields
\eqref{eq:coefficient-bound}.  For $|z|\le A$ and $R=2(A+1)$, the coefficient
bound gives
\begin{equation}\label{eq:egf-tail}
 \sum_{n>N}\frac{B_n}{n!}|z|^n
 \le 2\exp(e^R-1)\,2^{-(N+1)}.
\end{equation}
An explicit integer upper bound for the large constant is available by
$e<3$: for integral $R$, take
$\exp(e^R-1)<3^{3^R}$.  Hence choosing
$N=s+2\cdot3^R+4$ suffices.  For rational $R$, first increase it to an integer.
The finite coefficient identity, the exponential tails, and
\eqref{eq:egf-tail} now justify the analytic identity.  They also justify all
termwise differentiations on smaller bounded disks, by inserting the extra
polynomial factors $n(n-1)\cdots$ in the same geometric tail argument.
\end{proof}

The complete Bell polynomial can be defined without analytic vocabulary by
\[
 C_n(x_1,\ldots,x_n)
 =\sum_{\pi\text{ a partition of }[n]}
       \prod_{C\in\pi}x_{|C|}.
\]
Grouping by the number $m_j$ of size-$j$ blocks proves the finite formula
\begin{equation}\label{eq:block-multiplicity}
 C_n(x_1,\ldots,x_n)
 =\sum_{\substack{m_1,\ldots,m_n\ge0\\\sum j m_j=n}}
   \frac{n!}{\prod_{j=1}^n m_j!(j!)^{m_j}}
   \prod_{j=1}^n x_j^{m_j}.
\end{equation}
Indeed, order all $n$ labels, cut them into the indicated blocks, and quotient
by the $j!$ orders within each block and the $m_j!$ orders of equal-sized blocks.
Thus $C_n(1,\ldots,1)=B_n$, so the sum of its coefficients is $B_n$.
Expanding a finite exponential also proves
\[
 \sum_{n\ge0}C_n(x_1,\ldots,x_n)\frac{z^n}{n!}
 =\exp\left(\sum_{j\ge1}x_j\frac{z^j}{j!}\right)
\]
coefficientwise.  In the usual moment--cumulant language, all the cumulants of
the specified Poisson weights are $1$.  This is exactly the preceding formal
identity specialized to $x_j=1$.

\subsection{Hypergeometric notation, without extra functions}

For $n\ge1$, the hypergeometric statement on MathWorld is
\[
 B_n=e^{-1}\,{}_{n-1}F_{n-1}
     (\underbrace{2,\ldots,2}_{n-1};
      \underbrace{1,\ldots,1}_{n-1};1).
\]
Define the expression on the right by its factorial series.  Since the rising
factorials satisfy $(2)^{\overline j}/(1)^{\overline j}=j+1$, its $j$th term is
$(j+1)^{n-1}/j!=(j+1)^n/(j+1)!$.  Changing the index in
Theorem~\ref{thm:dobinski} proves the formula with the same effective tail.
For $n=1$ this includes ${}_0F_0(1)=e$.

\subsection{The ordinary series and its negative-axis substitute}

The ordinary generating series $\sum B_nx^n$ is formal.  It is not a
convergent series away from $0$.  Nevertheless, on the negative real axis
there is a useful function having that series as its asymptotic expansion.

\begin{theorem}[A Stieltjes-type representation and its full remainder]
For $x\le0$, define the effective name
\[
 G(x)=e^{-1}\sum_{k=0}^{\infty}\frac1{(1-kx)k!}.
\]
It is uniformly convergent on this half-line, with tail at most
$2/(M+1)!$ after $k=M$.  For every $N\ge0$,
\begin{equation}\label{eq:ordinary-asymptotic}
 G(x)=\sum_{j=0}^N B_jx^j+\mathcal R_N(x),\qquad
 |\mathcal R_N(x)|\le B_{N+1}|x|^{N+1}.
\end{equation}
The remainder has sign $(-1)^{N+1}$ for $x<0$.
\end{theorem}
\begin{proof}
Use the exact finite identity
\[
 \frac1{1-kx}=\sum_{j=0}^N(kx)^j+
                    \frac{(kx)^{N+1}}{1-kx}.
\]
Multiply by $e^{-1}/k!$, sum first over a finite range of $k$, and apply
Theorem~\ref{thm:dobinski}.  For $x\le0$ the denominator is at least $1$,
so the absolute value of the last sum is at most $B_{N+1}|x|^{N+1}$.
The factorial tail proves existence of the name independently of the
asymptotic expansion.  Each summand in the last sum has the claimed sign.
\end{proof}

For $x<0$ put $a=-1/x>0$.  Then
$1/(1-kx)=a/(a+k)=(a)^{\overline k}/(a+1)^{\overline k}$.
Consequently the confluent hypergeometric representation is exactly
\[
 G(x)=e^{-1}\,{}_1F_1(-1/x;1-1/x;1),
\]
with the defining series and the preceding tail bound.  We do not use the
incomplete-gamma expression from MathWorld, whose written form involves
complex branch conventions.

The other convergent series displayed there also follows by finite algebra:
\begin{equation}\label{eq:ordinary-rising-series}
 G(x)=\sum_{k=0}^{\infty}\frac{(-1)^k}{(a+1)^{\overline k}}
     =\sum_{k=0}^{\infty}
       \frac{x^k}{(-x)^k((x-1)/x)^{\overline k}}
       \qquad(x<0,\ a=-1/x).
\end{equation}
Here the rising factorial with index zero is one. To prove the identity,
define the finite rational sum
\[
 H_n(a)=\sum_{j=0}^{n}(-1)^j\binom nj\frac{a}{a+j}.
\]
Pascal's identity gives
\[
 H_{n+1}(a)=H_n(a)-\frac{a}{a+1}H_n(a+1).
\]
Starting with $H_0(a)=1$, induction therefore gives
\begin{equation}\label{eq:finite-beta-sum}
 H_n(a)=\frac{n!}{(a+1)^{\overline n}}:
\end{equation}
indeed, the difference of the two proposed terms has numerator
$n!((a+n+1)-a)=(n+1)!$ over
$(a+1)\cdots(a+n+1)$. For quantifier-free induction, fix the final index
$N$ and retain the finite table of assertions for arguments
$a,a+1,\ldots,a+N-r$ at stage $r$; each next entry uses two entries already
in that table. All denominators are positive and may be cleared.

The coefficient of $z^n$ in the finite product
$e^{-z}\,{}_1F_1(a;a+1;z)$ is
\[
 \frac1{n!}\sum_{j=0}^{n}(-1)^{n-j}\binom nj\frac{a}{a+j}
   =\frac{(-1)^n}{(a+1)^{\overline n}},
\]
by \eqref{eq:finite-beta-sum}. Since $0<a/(a+j)\leq1$ and
$(a+1)^{\overline n}\geq n!$, all three series and their product have
factorial tail bounds on every fixed bounded disk. Thus the finite coefficient
calculation passes to their effective names. Setting $z=1$ proves
\eqref{eq:ordinary-rising-series}, with absolute tail at most
$2/(M+1)!$ after index $M$. This adds no special-function assumption.

To see directly that the ordinary series has radius zero, fix a positive
rational $x$ and choose an integer $k$ with $kx\ge2$.  Dobi\'nski gives
\[
 B_nx^n\ge\frac{(kx)^n}{e\,k!}\ge\frac{2^n}{3k!}.
\]
An explicit $n$ makes this greater than any prescribed rational number.
The terms therefore do not tend to zero.  The same absolute-value argument
applies at nonzero complex $x$ with an explicit positive rational lower bound
for $|x|$.  This is also an elementary proof that Bell numbers grow faster
than every fixed exponential, independent of the sharper growth theorems.

\subsection{Fourier coefficient integrals and Ces\`aro's formula}

\begin{theorem}[Circular coefficient extraction]\label{thm:fourier-coefficient}
For $r>0$ and $n\ge0$,
\begin{equation}\label{eq:circular-extraction}
 B_n=\frac{n!}{2\pi r^n}
  \int_{-\pi}^{\pi}\exp(e^{r e^{i\theta}}-1)e^{-in\theta}\,d\theta.
\end{equation}
Equivalently, on the positively oriented circle $|z|=r$,
\[
 B_n=\frac{n!}{2\pi i e}\int_{|z|=r}
                    \frac{\exp(e^z)}{z^{n+1}}\,dz.
\]
For $n\ge1$ the real integral representation is
\begin{align}\label{eq:cesaro}
 B_n=\frac{2n!}{\pi e}\int_0^\pi
  &\exp\!\bigl(e^{\cos\theta}\cos(\sin\theta)\bigr)\notag\\[-3pt]
  &\quad\cdot\sin\!\bigl(e^{\cos\theta}\sin(\sin\theta)\bigr)
    \sin(n\theta)\,d\theta.
\end{align}
\end{theorem}
\begin{proof}
Insert the uniformly convergent series \eqref{eq:egf} on the circle.  The
finite orthogonality identity
$(2\pi)^{-1}\int_{-\pi}^{\pi}e^{i(m-n)\theta}\,d\theta=[m=n]$
extracts the coefficient, and \eqref{eq:egf-tail} bounds the discarded part.
Parametrizing the circle by $z=r e^{i\theta}$ proves the contour notation.
No general Cauchy theorem is required.

At $r=1$ the imaginary part of $\exp(e^{e^{i\theta}})$ is
$e\sum_{m\ge1}B_m\sin(m\theta)/m!$.  Multiplication by $\sin(n\theta)$
and the identity
$\int_0^\pi\sin(m\theta)\sin(n\theta)\,d\theta=(\pi/2)[m=n]$
give the result.  Finally expand the two nested complex exponentials into
their sine and cosine parts to obtain \eqref{eq:cesaro}.  This agrees with
the corrected formula in Callan~\cite{callan-cesaro}; the factor $n!$ is essential.
All integral identities here are identities of the effective Riemann-sum
names in Section~\ref{subsec:names}.
\end{proof}
